La sonda Voyager 1, llançada el 1977, va ser la primera nau en sortir del sistema solar després d'interac\-cionar amb Júpiter. El llançament va propulsar la seva massa de 720~kg a 10,0~km/s respecte la Terra en la mateixa direcció i sentit que la velocitat de la Terra respecte el Sol.

\begin{apartat}{1,25}
Considerant sempre òrbites planetàries circulars respecte el Sol, calcula el moment angular i la velocitat angular del Voyager 1 just després del llançament, encara a prop de l'òrbita terrestre. Calcula amb quina velocitat arribà el Voyager 1 a les proximitats de Júpiter i la seva velocitat angular en aquest punt.
\end{apartat}

\begin{apartat}{1,25}
A partir de l'expressió general de l'energia mecànica obtingueu l'expressió de la velocitat d'escapament del Sol per una nau a una distància $d_{nau-sol}$. Calculeu quina és aquesta velocitat per una nau a prop de Júpiter i decidiu si la Voyager 1 podia escapar-se del Sol abans d'interaccionar amb Júpiter. Calculeu l'energia mínima proporcionada per Júpiter a la nau per que aquesta s'hagi pogut escapar del sistema solar.
\end{apartat}

\dades{%
  $G = 6,67\cdot 10^{-11}\ \mathrm{N\,m^2\,kg^{-2}}$\\
  Massa del Sol $M_\mathrm{S} = 1,989\cdot 10^{30}\ \mathrm{kg}$\\
  Distància del Sol a la Terra $= 150\cdot 10^{9}$~m\\
  Distància del Sol a Júpiter $= 740\cdot 10^{9}$~m}
